\chapter{Performance, profiling, and the budget}
\label{ch:perf}

\begin{aside}
A TUI feels good when frames land in time. A 60Hz cycle gives
16.67ms; a 30Hz cycle 33.3ms. Most of that budget is the
terminal's own latency. \maya{}'s work fits in microseconds, but
only if you stay in the lanes the architecture invites. This
chapter is the budget: where the cycles go, how to measure,
how to fix what you find.
\end{aside}

\section{The frame budget}

A modern terminal commits a frame when:

\begin{enumerate}[leftmargin=*]
  \item Your program writes bytes to stdout.
  \item The terminal's render loop wakes (typically every
    frame).
  \item The GPU draws.
\end{enumerate}

The terminal contributes most of the latency: 5--15 ms is
normal. \maya{}'s pipeline runs in well under 1 ms. If your
frame total exceeds 16 ms, the work is almost certainly yours
(in \code{update} or \code{view}), not the library's.

\subsection*{Where the budget breaks down}

For a typical UI:

\begin{itemize}[leftmargin=*]
  \item \code{view}: 100--500 \textmu s.
  \item Layout: 20--100 \textmu s.
  \item Paint: 5--50 \textmu s.
  \item Diff: 1--10 \textmu s.
  \item Serialise: 5--50 \textmu s.
  \item Write to stdout: 50--200 \textmu s.
  \item Terminal latency: 5--15 ms.
\end{itemize}

Your work is the 100--500 \textmu s in \code{view}. The rest is
fixed cost.

\begin{idea}
The biggest single optimisation for most apps is moving
expensive computation out of \code{view} and into \code{update}
(cached in the model). \code{view} runs every frame; \code{update}
runs once per relevant message.
\end{idea}

\section{Profiling the right way}

Three measurement layers, from coarse to fine.

\subsection*{The pipeline stats}

\maya{}'s pipeline exposes per-frame timings via the \code{Stats}
struct. Enable the overlay:

\begin{lstlisting}
maya::config().show_render_stats = true;
\end{lstlisting}

A small widget appears in a corner with \code{view\_time},
\code{layout\_time}, \code{diff\_time}, \code{serialize\_time},
and frame total. Watch which number changes when you do
something.

\subsection*{Sampling profiler}

For deeper analysis, run under \code{perf} (Linux) or
\code{Instruments} (macOS):

\begin{lstlisting}
perf record -g ./my_maya_app
perf report
\end{lstlisting}

The flamegraph reveals which functions consume CPU. \code{view}
calls dominate; below them, the cost lives in your business
logic or in \code{element::build}.

\subsection*{Per-frame log}

For occasional spikes, log per-frame timings to a file:

\begin{lstlisting}
maya::config().log_frame_times = \"/tmp/frames.log\";
\end{lstlisting}

The log writes one line per frame with the timings. Grep for
outliers; correlate with what the user was doing.

\section{Optimisation playbook}

When you find a slow frame, the playbook in order:

\subsection*{1. Move work to update}

If \code{view} parses a string, formats a number, or walks a
big list, do it once in \code{update} and cache the result:

\begin{lstlisting}
// Bad: parses every frame.
Element view(const Model& m) {
    auto parsed = parse_huge_thing(m.input);
    return render(parsed);
}

// Good: cached in the model.
struct Model {
    std::string input;
    ParsedThing cached;  // refreshed in update when input changes
};
\end{lstlisting}

\subsection*{2. Virtualise long lists}

If \code{view} renders thousands of items, only render the
visible slice (Chapter~\ref{ch:scroll}). The layout pass scales
with the rendered count, not the source size.

\subsection*{3. Use compile-time DSL where possible}

A static structure built at compile time has fewer allocations
than the same structure built at runtime. For panels with fixed
shape, the DSL pays for itself.

\subsection*{4. Reuse buffers}

If a widget allocates a large vector every frame, capture the
vector in the widget's state and reset (\code{vec.clear()})
instead of reconstructing. Capacity persists; allocations
disappear.

\subsection*{5. Reduce allocations in update}

Update is hot. Every \code{std::string} concatenation allocates;
every \code{std::vector} push that exceeds capacity allocates.
Profile, then move construction outside hot loops.

\subsection*{6. Cache derived state}

Computed values that depend on slow inputs deserve a cache.
\code{Signal} + \code{Computed} (Chapter~\ref{ch:reactivity-preview})
gives you automatic recomputation only on change.

\subsection*{7. Tune the cache fingerprint}

If your tree is huge but mostly stable, the fingerprint walk is
the bottleneck. Two options: split the tree into smaller
sub-trees (each fingerprinted independently) or mark stable
sub-trees as \emph{frozen} (skipped in the walk).

\section{Memory budget}

A \maya{} app's resident memory:

\begin{itemize}[leftmargin=*]
  \item Model: usually small (KB). Watch for large vectors
    that grow unbounded (logs, history).
  \item Element tree: KB to low MB. Freshly allocated each
    frame; garbage-collected by RAII at frame end.
  \item Canvas: width $\times$ height $\times$ cell\_size. For
    80x24, that is 1920 cells $\times$ 8--16 bytes = a few KB
    each (current + previous = two copies).
  \item Style pool: a few KB.
  \item Renderer buffers: KB.
\end{itemize}

A small \maya{} app sits at 5--20 MB resident, dominated by the
binary itself. If you see hundreds of MB, look for unbounded
vectors in the model.

\subsection*{Allocations per frame}

A well-tuned app makes a handful of allocations per frame:

\begin{itemize}[leftmargin=*]
  \item Element tree allocations: a vector per box,
    \code{string\_view}s for literals.
  \item Layout vector resize (cap reused).
  \item Render buffer resize (cap reused).
\end{itemize}

\code{strace -c -e brk,mmap} on a running app shows the rate.
If it climbs into the thousands per second, hunt the
allocation sources.

\section{Latency vs throughput}

The pipeline optimises for \emph{latency} (frame-to-frame
time), not throughput (frames per second). The trade-offs:

\begin{itemize}[leftmargin=*]
  \item Low-priority diff: skip if last frame was less than
    16ms ago. Reduces CPU on busy apps; adds latency.
  \item Triple buffering: render frame N+2 while frame N is on
    the wire. Adds latency; increases throughput.
\end{itemize}

\maya{}'s defaults favour latency. If you have a soft-realtime
requirement (audio visualiser, game), tune the renderer
explicitly.

\section{Common slowdowns}

\subsection*{Reading large files in view}

The classic. \code{view} opens a file, reads it, formats it,
returns the element. Every frame. The file is re-read 60 times
a second. Cache the contents in the model and refresh only when
the file changes (use \code{Cmd::task} with a file-watch).

\subsection*{Heavy regex in view}

Regex compilation and matching are slow. Cache the compiled
pattern; cache the match results.

\subsection*{Recursive view functions}

A \code{view} that calls itself deeply on a recursive data
structure can blow stack and time. Convert to iterative;
materialise levels separately.

\subsection*{Mouse events firing on every move}

\code{OnMouse} fires for every motion event, which can be
thousands per second on a fast trackpad. Throttle the message
delivery (\code{Sub::throttle(30ms, OnMouse)}) or filter at the
handler.

\subsection*{Style pool exhaustion}

The style pool interns styles to small IDs. If you generate
thousands of distinct styles (e.g. one per character of a
gradient), the pool grows and dispatches slow down. Reuse a
few common styles where possible.

\section{When to give up tuning}

If you have:

\begin{itemize}[leftmargin=*]
  \item Frame time under 5 ms.
  \item No visible jank to the user.
  \item No CPU usage during idle.
\end{itemize}

\ldots{} you are done. Further optimisation is anti-economic.

\maya{}'s architecture lets you stop tuning earlier than other
frameworks because the defaults are aggressive: cache
fingerprints, SIMD diff, batched serialise. Profile only the
remaining hot spots in your own code.

\begin{aside}
Donald Knuth's ``premature optimisation is the root of all evil''
is often misquoted to argue against any optimisation. The full
quote ends ``yet we should not pass up our opportunities in that
critical 3\%''. Identify the 3\%; tune it; leave the rest.
\end{aside}

\section{Benchmarking widgets}

For widgets you ship to others, include a benchmark:

\begin{lstlisting}
BENCHMARK(\"render textarea with 1000 lines\") {
    Model m = make_model_with_lines(1000);
    return view(m);
}
\end{lstlisting}

The benchmark suite measures across compilers, terminal sizes,
and content. Track regressions over time.

\section*{Workshop}

\begin{enumerate}[leftmargin=*]
  \item Enable \code{show\_render\_stats} on a small app.
    Note the per-stage timings. Identify which stage dominates.
  \item Add a deliberately slow function to \code{view} (e.g.
    \code{std::this\_thread::sleep\_for(10ms)}). Watch the
    frame time rise.
  \item Profile a real app under \code{perf}. Identify the top
    three functions; fix the slowest one.
\end{enumerate}

\begin{recap}
The pipeline runs in microseconds; user work is the variable
cost. Move expensive computation from \code{view} to
\code{update} with caching. Virtualise long lists. Reuse
buffers. Avoid recomputing what has not changed. Memory budget
is small. The defaults are aggressive; tune only the residual
hot spots.
\end{recap}
